% brainstorm:

% * revisão das teorias de contornos
% * parte da análise que fiz do opus 5 mov 3 de Webern
% * exemplos de usos de operações na peça do mestrado e experimentos
% * desenvolvimento do goiaba (para que serve e um exemplo da saída)

\section{Introduction}
\label{sec:introduction}

Contours can be understood as the shape or format of an object. In
Music, contour can be associated to elements like pitch, density,
rhythm, intensity, and represents a parameter in function of another,
as pitch in function of time. For instance, Beethoven's Fifth Symphony
main motive and its pitch contour are represented respectively in
figures \ref{fig:5a-sinfonia-motivo} and \ref{fig:c-3120}.

Absolute value of contour elements are ignored, and only the high-low
relationship between these elements is regarded. For this reason, for
example, figures \ref{fig:5a-sinfonia-motivo} and \ref{fig:ly-3120}
have the same contour, represented in figure \ref{fig:c-3120}. In
these figures the first pitch is highest, the second is lower than the
first, the third is between first and second, and the last is the
lowest.

The study of Contour is important because, like motives and pitch
sets, contours can help to give coherence to a musical piece. They
represent structural devices that can be combined through operations
like inversion and retrogradation, and can be aproached by analytical
or compositional points of view.

%% melhorar justificativa
Despite the possible coherence that contours can give, and the
operations provided by contour theories (see section
\ref{sec:contour-theories}), systematic studies of contour operations
and combinations usage in musical composition are scarce. So more
experimentation and studies about these operations usage in the field
of composition are required.

In this paper we present partial results of our research about contour
applications in musical composition. We are developing a contour
processor software\footnote{Additional information and link were
  omitted to keep anonymity. It will be included in final version.}
that receives as input symbollic contours and outputs symbolic and
graphical representations of contours and operations. The first author
of this paper, in his master's research\footnote{The reference was
  also omitted to keep anonymity.}, composed an eleven minutes quintet
based in contour theories operations using the mentioned software (see
section \ref{sec:contour-composition}).

\section{Contour theories}
\label{sec:contour-theories}

Many authors
\cite{friedmann85:methodology,friedmann87:response,morris87:composition,morris93:directions,marvin.ea87:relating,marvin88:generalized,marvin.ea95:generalization,polansky.ea92:possible,quinn97:fuzzy,clifford95:contour,beard03:contour}
have developed theories to organize in a systematic way the knowledge
about contour. These theories were developed primarily as analytic
techniques for non-tonal compositions \cite{beard03:contour}. These
theories provide arrays, matrices and many operations to help the
comparison of contours, like inversion, translation, comparison of
elements of a contour, comparison of simalarity of contours,
comparison matrix and contour interval array.

Music analysis under the point of view of contour have been effective
in tonal and non-tonal music. There are successful analysis of pieces
by A. Schoenberg \cite{friedmann85:methodology}, A. Webern
\cite{clifford95:contour}, L. Dallapicolla
\cite{marvin88:generalized}, and W. A. Mozart \cite{beard03:contour}
pieces. Clifford \cite{clifford95:contour}, for example, says contour
has a significant role in Webern pre-serial music structure.

\section{The usage of contour in Composition}
\label{sec:contour-composition}

The cited eleven minutes quintet, composed during research, is based
on combinations of contour operations associated with parameters such
as pitch, tempo, density and texture. This quintet is based on a six
notes motive $\alpha$ (fig. \ref{fig:motivo-alfa})---based on
octatonic scale---and its derived contour P(5 3 4 1 2 0)
(fig. \ref{fig:c-534120}). This P contour, its subsets and
operations---retrogradation, inversion, rotation,
interpollation---were the only one used in the piece.

For instance, the tempo in this quintet---82, 66, 120, 108 and
112---have relations based in A(1 0 4 2 3) contour, a subset of P with
five elements. The number of instruments in piece first section, N(1 3
2 5 4), also are based in a P subset contour. The piece begins with a
solo, then a trio, a duo, a quintet, and finally a quartet.

In the quintet, contour operations were also combined to produce new
material. Rotation were combinated with retrogradation, intervallic
expansion with interpolation, rotation with intervallic expansion, and
so on.

For instance, in a \eng{fugato}, in the quintet, subject and
countersubject are based on rotation and retrogradation
operations. The subject has original contour P(5 3 4 1 2 0) and factor
3 rotation (fig. \ref{fig:sujeito-fugato}). The contoursubject has
original contour retrogradation repeated three times with different
rotation factors---5, 4 and 3 (fig. \ref{fig:contra-sujeito-fugato}).

Rotation and intervallic expansion occur like in figure
\ref{fig:notas-curtas-madeiras}. In flute the contour has original
form, and in clarinet there is a factor 2 rotation and an
expansion. In oboe there is a factor 3 rotation and contour deviation
justified by the gesture. In bassoon there is a factor 3 contour
rotation and retrogradation, and intervallic expansion.

\section{Conclusions}
\label{sec:conclusions}

Contour represented a significant structural device in cited quintet
compositional process. The majority structures of the piece were
composed based on contour operations. For this reason we believe
contour is as important to composition as to analysis.

The cited software has an important role in our research. Contour can
be easily recognized from graphic representation (as in figures
\ref{fig:rot-0}, \ref{fig:rot-2}, \ref{fig:rot-3}, and
\ref{fig:rot-3-retr}). Also, contour operations requires precise
mathematical handling. This software automates both tasks avoiding
human errors and wasting time. In the future the software will receive
as input scores in the Lilypond\footnote{\url{http://lilypond.org/}}
format and will have a graphical interface.

The next step in our research is to develop the software to provide a
better interaction with the user, and to compose short experiments to
test a larger operations corpus, and multiple layers contours, like
simultaneous association with pitch, duration, timbre and dynamics.

\break
%%% concentra figuras em um só lugar

\begin{figure}[!p]
  \centering
  \subfloat[Main motive]{
    \includegraphics{5a-sinfonia}
    \label{fig:5a-sinfonia-motivo}
  }
  \subfloat[Contour (3 1 2 0)]{
    \includegraphics{c-3120}
    \label{fig:c-3120}
  }

  \subfloat[Another (3 1 2 0) melody]{
    \includegraphics{ly-3120-qualquer}
    \label{fig:ly-3120}
  }
  \caption{Fifth Symphony main motive contour}
  \label{fig:5a-sinfonia}
\end{figure}

\begin{figure}[!p]
  \centering
  \subfloat[$\alpha$ motive]{
    \includegraphics{motivo-alfa}
    \label{fig:motivo-alfa}
  }
  \subfloat[P(5 3 4 1 2 0) contour]{
    \includegraphics{c-534120}
    \label{fig:c-534120}
  }
  \caption{Materials}
  \label{fig:materials}
\end{figure}

\begin{figure}
  \centering
  \subfloat[Subject]{
    \includegraphics[scale=3.2]{sujeito-fugato}
    \label{fig:sujeito-fugato}
  }

  \subfloat[Countersubject]{
    \includegraphics[scale=3.2]{contra-sujeito-fugato}
    \label{fig:contra-sujeito-fugato}
  }
  \caption{Structural elements of \eng{fugato}}
  \label{fig:elementos-fugato}
\end{figure}

\begin{figure}[!p]
  \centering
  \subfloat[4 fragments composed by rotation and expansion]{
    \includegraphics[scale=1]{notas-curtas-madeiras}
    \label{fig:notas-curtas-madeiras}
  }

  \subfloat[original]{
    \includegraphics[scale=.75]{c-534120}
    \label{fig:rot-0}
  }
  \subfloat[rot 2]{
    \includegraphics[scale=.75]{c-412053}
    \label{fig:rot-2}
  }
  \subfloat[rot 3]{
    \includegraphics[scale=.75]{c-120534}
    \label{fig:rot-3}
  }
  \subfloat[retr(rot 3)]{
    \includegraphics[scale=.75]{c-435021}
    \label{fig:rot-3-retr}
  }
  \caption{Rotation and expansion}
  \label{fig:rotacao-expansao}
\end{figure}

%%% Local Variables: 
%%% mode: latex
%%% TeX-master: "contour-composition"
%%% End: 